\section{Bihomogeneous determinant cuts}
\label{sec:determinant-cuts}

We now reduce the dimension of the algebraic sets containing the points
in each paired box of Lemma~\ref{lem:paired-boxes}.  The first
determinant gives a hypersurface section.  When one component of that
section is defined using only one coordinate block, a second determinant
reduces its dimension again.  The estimates for this second step must be
uniform in the equation of the moving component.
The Taylor-determinant construction belongs to the real determinant
method of Bombieri--Pila and its approximating form developed by
Heath-Brown \cite{BombieriPila1989,HeathBrown2009Sums}; the
number-field application follows the arrangement in
\cite[Section~4]{Reuss2015PowerFree}. We supply the complete estimates,
including the second cut on pure components.

\subsection{The bicone and its Hilbert polynomial}

Fix the distinct roots \(\theta_1,\ldots,\theta_d\), and put \(k=d-2\).
In \(\C^{2d}\), with coordinate blocks \(a=(a_i)\), \(b=(b_i)\), let
\(\mathcal B\) be the scheme defined by the condition
\[
 (a_i^kb_i)_{i=1}^d
 \in \operatorname{span}_{\C}\{(1)_i,(\theta_i)_i\}.
\]
Equivalently, impose \(d-2\) independent linear relations on these
products.  Each relation has bidegree \((k,1)\).  Write the products as
\[
 a_i^kb_i=N_0-\theta_iL_0;
\]
on \(\mathcal B\), the functions \(N_0,L_0\) are fixed linear
combinations of the products.  The corresponding subvariety of
\(\mathbb P^{d-1}\times\mathbb P^{d-1}\) will be denoted by
\(\mathcal V\).

\begin{lemma}\label{lem:bicone}
The defining ideal of \(\mathcal B\) is prime and is a complete
intersection of \(d-2\) equations of bidegree \((k,1)\).
The dimension of \(\mathcal V\) is \(d\).  If \(A,B\) tend to infinity
with fixed positive rational ratio \(B/A\), its bihomogeneous coordinate
space in degree \((A,B)\) has dimension
\[
 R(A,B)+O(A^{d-1}),
\]
where
\begin{equation}
 R(A,B)=\frac1{d!}\sum_{j=1}^{d-1}
 \binom dj\binom{d-2}{j-1}k^{j-1}A^{d-j}B^j.
 \label{eq:hilbert-polynomial}
\end{equation}
\end{lemma}

\begin{proof}
Where all \(a_i\ne0\), the coordinates \(a_i,N_0,L_0\) are free and
determine the \(b_i\).  This is an integral smooth open of dimension
\(d+2\).  If some \(a_i=0\) and the product vector is nonzero, then
\(N_0=\theta_iL_0\), \(L_0\ne0\), and every other \(a_j\ne0\), since
the roots are distinct.  This stratum has the \(d-1\) coordinates
\(a_j\), together with \(L_0,b_i\), and dimension \(d+1\).
If the product vector is zero, the equations are
\(a_i^kb_i=0\) for all \(i\).  Their zero set is a union of coordinate
\(d\)-planes.

These counts give codimension \(d-2\) and a unique top-dimensional
component meeting the first open.  Here and below we use the following
complete-intersection criterion.  In a polynomial ring, homogeneous
equations of the expected codimension form a regular sequence; their
quotient is equidimensional and Cohen--Macaulay, with no embedded
components.  If there
is one top-dimensional component and the scheme is reduced at its
generic point, the quotient is reduced and hence a domain.
For completeness, localize at the homogeneous maximal ideal.  The
ambient regular local ring is Cohen--Macaulay, so the expected
codimension gives a regular sequence; the positive grading detects
regularity globally.  The quotient is Cohen--Macaulay and
equidimensional by the local dimension formula
\cite[Tags 00NQ, 02JN, 00NB, and 00NA]{stacks-project}.
It therefore has no embedded associated primes.  Generic reducedness
on its unique minimal component then implies reducedness
\cite[Tags 031Q and 031R]{stacks-project}.
This reasoning also applies to positive weighted homogeneous equations.
In the present case the indicated smooth open supplies generic
reducedness, proving primeness.  Passing to the two projective factors
subtracts two dimensions.

The Hilbert series of the regular sequence is
\[
 \frac{(1-z_1^kz_2)^{d-2}}{(1-z_1)^d(1-z_2)^d}.
\]
Thus the dimension in bidegree \((A,B)\) is the \((d-2)\)-fold
finite difference, with step \((k,1)\), of
\(\binom{A+d-1}{d-1}\binom{B+d-1}{d-1}\).
Its leading term is
\[
 (k\partial_A+\partial_B)^{d-2}
 \frac{A^{d-1}B^{d-1}}{((d-1)!)^2},
\]
which expands to \eqref{eq:hilbert-polynomial}.  The remaining terms
have total degree at most \(d-1\).
\end{proof}

Define the one-variable polynomials
\[
 H(t)=d!R(1,t),\qquad
 H_a(t)=(d-1)!(\partial_A R)(1,t),\qquad
 H_b(t)=(d-1)!(\partial_B R)(1,t).
\]
They measure the spaces of determinant columns on \(\mathcal V\)
and on a pure hypersurface section, respectively.

\begin{proposition}\label{prop:archimedean-cuts}
Fix positive upper exponents $\eta_+,b_+$.  Consider a paired box
as in Lemma~\ref{lem:paired-boxes}, with
integer coordinates satisfying
\[
 |x_i|\ll X^{\eta_+/d},\qquad |y_i|\ll X^{b_+/d}.
\]
Assume that its normalized chart variables have width \(O(X^{-w})\),
that \(z=1/n=O(X^{-1})\), and that the forward and reverse analytic
chart functions have the uniform derivative bounds of that lemma.
Let \(t_*,w>0\) be fixed rational numbers satisfying
\begin{equation}
 \begin{split}
 w^{d-1}H(t_*)&>
 \left(\frac{(d+1)(\eta_++b_+t_*)}{d^2}\right)^d,\\
 w^{d-2}\min\{H_a(t_*),H_b(t_*)\}&>
 \left(\frac{\eta_++b_+t_*}{d-1}\right)^{d-1}.
 \end{split}
 \label{eq:cut-conditions}
\end{equation}
For sufficiently large \(X\), the points in this box lie on a bounded
number of bounded-degree algebraic sets of the following two types:
\begin{enumerate}
\item a mixed hypersurface section of \(\mathcal V\), whose bicone is
a prime complete intersection with one additional equation of
bidegree \((\lambda,s_0)\), where \(\lambda,s_0\geq1\);
\item a set of biprojective dimension at most \(d-2\).
\end{enumerate}
On returning to the affine equations \(a_i^kb_i=n-\theta_i\), the
corresponding dimensions are \(d\) and at most \(d-1\).
All number and degree bounds are uniform over the boxes and \(X\).
\end{proposition}

\subsection{A uniform Taylor determinant estimate}

We first describe the analytic estimate used in both cuts.
Let \(\ell\) column functions have uniformly bounded derivatives on
the chart neighborhoods, and let the small variables be
\(u_1,\ldots,u_{d-1},z\), of weights \(w,\ldots,w,1\).
For a fixed finite set of evaluation points, suppose that the rank
of the evaluation columns of monomials of weight at most \(T\) is
at most
\[
 aT^r+O(T^{r-1}+1),\qquad a>0.
\]
If \(\ell\) such monomial columns are independent, and their
weights in increasing order are \(\lambda_j\), then
\(j\leq a\lambda_j^r+O(\lambda_j^{r-1}+1)\).  Inverting and summing
gives
\begin{equation}
 \sum_{j=1}^{\ell}\lambda_j
 \geq \frac{r}{r+1}a^{-1/r}\ell^{1+1/r}-O(\ell).
 \label{eq:taylor-rank-weights}
\end{equation}
The constants are uniform whenever the rank estimate is uniform.

Expand each analytic column at the box base in the \emph{original}
monomials in \(u,z\), to a fixed total order \(K\).
A polynomial summand in the multilinear expansion of the determinant
vanishes if its monomial evaluation columns are dependent.
Otherwise \eqref{eq:taylor-rank-weights} bounds its total weight, and
hence its absolute value by a constant times \(X\) to the negative
of that weight.  The Taylor coefficients are uniformly bounded.
The remainder in a column is
\(O(X^{-\min(w,1)(K+1)})\).
After fixing the column degrees and \(\ell\), choose \(K\) so large
that every determinant term containing a remainder has a stronger
bound than the required one.  There are then only boundedly many
terms.  Thus the normalized determinant is
\begin{equation}
 \ll X^{-S},\qquad
 S=\frac{r}{r+1}a^{-1/r}\ell^{1+1/r}-O(\ell).
 \label{eq:taylor-determinant-bound}
\end{equation}
This argument requires only exact dependence of evaluation columns,
not any bound on coefficients expressing that dependence.

\subsection{The first cut and its components}

Choose large integer bidegrees with \(B/A=t_*\), and choose a basis
of the degree-\((A,B)\) coordinate space of \(\mathcal V\) from raw
monomials in the integral coordinates \(x,y\).
By Lemma~\ref{lem:bicone}, its size is
\[
 \ell=\frac{H(t_*)}{d!}A^d+O(A^{d-1}).
\]
For any \(\ell\) counted points, the raw evaluation determinant is
an integer.  Divide each row by \(x_0^Ay_0^B\), using the nonzero
chart coordinates.  If the determinant is nonzero, its normalized
absolute value is at least
\begin{equation}
 c_{A,B}\,X^{-\ell A(\eta_++b_+t_*)/d},
 \qquad c_{A,B}>0.
 \label{eq:integer-determinant-lower}
\end{equation}
The constant can be chosen uniformly in the box.

Normalized monomials are analytic functions of the chart variables
and \(z\).  Without an extra relation among the small variables,
their monomial count is
\[
 \frac{T^d}{d!w^{d-1}}+O(T^{d-1}+1).
\]
Equations \eqref{eq:taylor-rank-weights}--\eqref{eq:taylor-determinant-bound}
therefore give
\[
 S=\frac d{d+1}(d!w^{d-1})^{1/d}\ell^{1+1/d}-O(\ell),
 \qquad
 \frac S\ell
 =\frac d{d+1}\bigl(w^{d-1}H(t_*)\bigr)^{1/d}A+O(1).
\]
The first inequality in \eqref{eq:cut-conditions} makes this exponent
larger than the one in \eqref{eq:integer-determinant-lower}, first
for sufficiently large fixed \(A,B\), then for sufficiently large
\(X\).  Every such determinant vanishes.  The evaluation map
therefore has rank less than \(\ell\), also when there are fewer
than \(\ell\) points.  Its kernel gives a bihomogeneous equation
vanishing on all the counted points but not identically on
\(\mathcal V\).

We must describe the actual components of this cut before using
their Hilbert functions.  On the open where all \(a_i\ne0\),
substitute \(b_i=(N_0-\theta_iL_0)/a_i^k\), clear denominators, and
factor.  Each relevant irreducible factor
\[
 F(a,N_0,L_0)=\sum_I a^If_I(N_0,L_0)
\]
is bihomogeneous, say of degrees \((l,s_0)\) in \(a\) and
\((N_0,L_0)\).  Indeed, the two scaling actions preserve each
irreducible factor up to scalar.  Factors \(a_i\) or \(L_0\) do
not meet the open of counted points and can be omitted.
Returning to the \((a,b)\)-coordinates replaces
\(N_0-\theta_iL_0\) by \(a_i^kb_i\), so the substituted terms may
acquire common powers of \(a_i\).  We extract these powers and then
verify that the remaining equation has no extra boundary component.
For each nonzero \(f_I\), set
\[
 \nu_{i,I}=\operatorname{ord}_{N_0-\theta_iL_0}f_I,\qquad
 e_i=\min_I(I_i+k\nu_{i,I}).
\]
Writing
\(f_I=\prod_i(N_0-\theta_iL_0)^{\nu_{i,I}}g_I\), the expression
\begin{equation}
 C(a,b)=\sum_I
 a^{I+k\nu_I-e}b^{\nu_I}g_I(N_0,L_0)
 \label{eq:saturated-bihomogeneous-cut}
\end{equation}
is a polynomial: every exponent of \(a\) is nonnegative.
It represents \(F/\prod_i a_i^{e_i}\) on the bicone, and has
bidegree
\[
 (\lambda,s_0),\qquad \lambda=l+ks_0-\sum_i e_i.
\]

This extension has no unwanted boundary component.  On \(a_i=0\)
with nonzero product vector, set \(L_0=1\).  Substitution for the
other \(b_j\) leaves, precisely when \(I_i+k\nu_{i,I}=e_i\),
a nonzero constant times
\[
 a_{-i}^{I_{-i}-e_{-i}}b_i^{\nu_{i,I}}.
\]
These Laurent monomials are distinct: their exponents determine
\(I_{-i}\) and \(\nu_{i,I}\), and then
\(I_i=e_i-k\nu_{i,I}\).  Thus the restriction of \(C\) is nonzero.
It reduces this boundary stratum to dimension at most \(d\);
the zero-product boundary already has dimension \(d\).
The open defined by \(F\) is integral and generically reduced of
dimension \(d+1\).  Since \(C\) is a nonzero element of the prime
bicone ring, adding it preserves the complete-intersection property.
The boundary calculation and the criterion in the proof of
Lemma~\ref{lem:bicone} now prove that the new ideal is prime.

The Hilbert series acquires the factor
\(1-z_1^\lambda z_2^{s_0}\).  Its leading term in bidegree \((A,B)\) is
\begin{equation}
 (\lambda\partial_A+s_0\partial_B)R(A,B),
 \label{eq:cut-hilbert-leading}
\end{equation}
with error \(O(A^{d-2})\) at fixed \(B/A\).
The nonnegative integers \(\lambda,s_0\) are not both zero.
If both are positive we retain this mixed component.
All degrees remain bounded, by the degree of the first determinant
equation and the fixed substitutions.

\subsection{The second cut on pure components}

Suppose instead that \(C\) uses only one block, with degree \(c_0>0\).
Use the \(s\)-chart for a pure \(a\) equation and the \(t\)-chart
for a pure \(b\) equation.  In the latter case, partition the counted
points according to the boundedly many reverse branches at the common
\(t\)-box base supplied by Lemma~\ref{lem:paired-boxes}, and apply the
following determinant argument to each subset.  Thus all rows of each
determinant use the same analytic chart map.  This only multiplies the
number of resulting sets by a bounded factor.
Dehomogenization and translation to the
box base give an exact polynomial relation
\[
 Q_0(u_1,\ldots,u_{d-1})=0,\qquad 1\leq\deg Q_0\leq c_0.
\]
A nonzero constant relation gives no points.  No bound on the
coefficients of \(Q_0\) is imposed.

Here is the needed coefficient-uniform rank estimate.  Put
\(r=d-1\) and \(e=\deg Q_0\).
Multiplication by \(Q_0\) injects polynomials of degree at most
\(q-e\) into polynomials of degree at most \(q\), and its image
evaluates to zero.  Therefore the evaluation rank of the latter
space is at most
\[
 D_e(q)=\binom{q+r}{r}-\binom{q-e+r}{r},
\]
with the second term zero when \(q<e\).
At each power \(z^j\), apply this bound with
\(q=\lfloor(T-j)/w\rfloor\).  Since
\(D_e(q)=e q^{r-1}/(r-1)!+O_{r,c_0}((q+1)^{r-2})\), summing gives
\[
 N(T)\leq
 \frac{c_0}{r!w^{r-1}}T^r+O_{r,c_0,w}(T^{r-1}+1).
\]
This bound involves no coefficients of the relation.
Inserting it into the Taylor argument above gives
\[
 S\geq\frac r{r+1}
 \left(\frac{r!w^{r-1}}{c_0}\right)^{1/r}
 \ell^{1+1/r}-O(\ell).
\]
In particular, we have never solved the moving equation for one
variable or reduced Taylor monomials by division by a coefficient.

Keep the first-cut degrees fixed, so that \(c_0\) ranges over a
bounded set of positive integers.  Choose new, sufficiently large
bidegrees \(A,B\), again with \(B/A=t_*\), for the second determinant.
By \eqref{eq:cut-hilbert-leading}, a basis of raw monomials on a
pure \(a\) component has size
\[
 \ell=\frac{c_0H_a(t_*)}{r!}A^r+O(A^{r-1}).
\]
Consequently
\[
 \frac S\ell\geq
 \frac r d\bigl(w^{r-1}H_a(t_*)\bigr)^{1/r}A-O(1).
\]
The factor \(c_0\) cancels.  The identical calculation for a pure
\(b\) component uses \(H_b\).  The second inequality in
\eqref{eq:cut-conditions} contradicts
\eqref{eq:integer-determinant-lower} in both cases.
The resulting equation is nonzero on the component, so it cuts
its biprojective dimension from \(d-1\) to at most \(d-2\).
The Hilbert series and errors depend only on bounded bidegrees,
and the analytic bounds come from the ambient chart maps.
Thus the second-cut degrees can be fixed uniformly for every
pure component.

Finally, on the counted affine locus \(L_0=1\), the map to
\(\mathcal V\) has one-dimensional scaling fibers:
\(a\mapsto ca\) requires \(b\mapsto c^{-k}b\).
Mixed components therefore have affine dimension \(d\), and
second-cut components have dimension at most \(d-1\).
In the mixed case the equation on the nonzero-\(a\) open is
\(F(a,n,1)=0\).  It is irreducible, since a factorization would
rehomogenize to one of \(F\), and its degree in \(n\) is exactly
\(s_0\), since \(L_0\) does not divide \(F\).
Since \(s_0\geq1\), its coefficient polynomials in \(a\) have no
common nonconstant factor; such a factor would contradict
irreducibility.  Their common zero set therefore has codimension
at least two in the \(a\)-space whenever it is nonempty.
Taking closures preserves these dimensions and bounded degrees.
This proves Proposition~\ref{prop:archimedean-cuts}.
